\section{Feature Engineering}
\label{sec:features}

\subsection{The 36 features, and why each is there}

Phase 1's 25 features carry forward unchanged: 6 calendar (month, weekday, hour,
weekend, business hours, the Black Friday window), 11 basket-and-money
(\code{log\_price}, \code{freight\_ratio}, \code{log\_weight}, volume, item and
seller counts, instalments), 4 geographic (\code{log\_distance},
\code{same\_state}, customer state, seller region), 2 for the promise and 2
catalogue. The design carries no exact duplicates, and \code{purchase\_year} is
absent because the split is chronological, so a linear term on it would
extrapolate into an unseen year. Phase 2 added \textbf{eight track-record
features}, taking the checkout set to 33 and the at-delivery set to 36: a
seller's mean lead time, late rate, low-review rate and order count; a route's
mean lead time; a category's low-review rate; and the platform's own lead time
and late rate. The at-delivery regime additionally sees three columns describing
how delivery actually went.

Each group has a reason in the literature of Section~\ref{sec:lit} and a
screen in the data. Geography and the promise are the backbone of delivery-time
models \citep{salari2022delivery, liu2021lastmile, cui2024promise}; basket,
money and calendar are the order-level controls of the logistics--rating
studies \citep{deshpande2023logistics}; track records are the standard
predictor in supply-chain delay models \citep{baryannis2019risk} in the
time-respecting form of target encoding \citep{miccibarreca2001target}. The
empirical screen was the correlation matrix of Figure~\ref{fig:correlations}:
lead time correlates 0.39 with distance, 0.38 with the promised window, 0.17
with freight and 0.07 with weight; promise and distance correlate 0.52 with
each other, so both are kept and the trees resolve the overlap.
Table~\ref{tab:features} in Appendix~\ref{app:features} lists all 36 features
with definition, rationale and the moment each becomes knowable.

\subsection{Track records, and the trap inside them}
\label{sec:history}

A seller's history is the strongest thing a marketplace knows about an order at
checkout, and the easiest place here to leak, because of a distinction that is
easy to miss: \emph{past} means the earlier order had \textbf{closed} before this
one was placed --- not merely that it was placed earlier. A parcel bought last
week that has not yet arrived has no lead time yet, so ordering the history by
\emph{purchase} time would let every order average in outcomes that had not
happened. That is Kaufman et al.'s \emph{illegitimate feature}
\citep{kaufman2012leakage} in a convincing disguise, and it would be invisible in
the metrics because it would simply look like a better model. Every aggregate is
therefore an as-of join against the moment the earlier outcome became observable
--- \code{order\_delivered\_customer\_date} for delivery,
\code{review\_creation\_date} for reviews --- with a strict boundary.

Two further guards. A mean over one previous order is noise rather than a rate,
so below a floor of five prior orders the value is undefined and a fallback chain
takes over (seller $\rightarrow$ route $\rightarrow$ seller state $\rightarrow$
platform), after which 0.95\% of orders lack any history value; and cold start
is information rather than error, since \code{hist\_seller\_n} carries how many
past orders the estimate rests on. The platform features use a \textbf{28-day
trailing window} rather than an expanding mean, because lead time halves across
the analysis window, and their range, 5.2 to 17.1 days, is a direct measure of
the drift the models must survive. Because a leak here would be invisible in
the metrics, the construction is checked rather than reviewed: the validation
suite rebuilds \code{hist\_seller\_lead\_days} from first principles for 200
random orders and requires exact agreement, and a second check confirms the
platform tracker \emph{lags} rather than anticipates (Appendix~\ref{app:repro}).
